\documentclass[tikz,border=6pt]{standalone}
\usepackage{peabra}
\begin{document}
\begin{tikzpicture}[x=1cm,y=1cm, a/.style={doboz=52mm, font=\footnotesize, minimum height=9mm}, b/.style={sotetdoboz=52mm, font=\footnotesize, minimum height=9mm}]
  \node[a] (ksz) at (0,4.6) {hatókör, környezet, kritériumok\\(mi a kockázatvállalási hajlandóság?)};
  \draw[peKek, line width=0.8pt, rounded corners=3pt, fill=peHalvany!40] (-3.15,0.35) rectangle (3.15,3.75);
  \node[font=\cimbetu\footnotesize, text=peKek, anchor=north west] at (-3.05,3.7) {KOCKÁZATÉRTÉKELÉS};
  \node[a] (az) at (0,2.85) {kockázatazonosítás\\(veszélyek, forgatókönyvek)};
  \node[a] (el) at (0,1.85) {kockázatelemzés\\(valószínűség, következmény)};
  \node[a] (er) at (0,0.95) {kockázatminősítés\\(összevetés a kritériumokkal)};
  \node[b] (ke) at (0,-0.6) {kockázatkezelés\\(csökkentés, elkerülés, megosztás, vállalás)};
  \draw[nyil] (ksz) -- (az); \draw[nyil] (az) -- (el); \draw[nyil] (el) -- (er); \draw[nyil] (er) -- (ke);
  \draw[nyil, draw=peKiemel, dash pattern=on 3pt off 2pt] (ke.east) -- ++(0.8,0) |- node[felirat, pos=0.25, right, text=peKiemel, align=left] {maradó\\kockázat\\újraértékelése} (el.east);
  \node[szurkedoboz=58mm, minimum height=9mm, rotate=90, font=\footnotesize] at (-4.0,2.0) {kommunikáció és konzultáció az érdekelt felekkel};
  \node[szurkedoboz=58mm, minimum height=9mm, rotate=90, font=\footnotesize] at (5.9,2.0) {nyomon követés és felülvizsgálat};
  \node[szurkedoboz=112mm, font=\footnotesize] at (0.95,-1.85) {dokumentálás és jelentés (kockázati nyilvántartás)};
\end{tikzpicture}
\end{document}
